\chapter{Equações de retas e sistemas lineares}\label{ch:g10:lines}

Uma reta no plano cartesiano é descrita por uma \omterm{def:g10:algebra:equation}{equação}, e duas retas
juntas formam um sistema de \omterm{def:g10:algebra:equation}{equações} cujas soluções são seus pontos de
\omterm{def:g10:numbers:interunion}{interseção}. Este capítulo conecta três linguagens --- a geométrica
(retas), a algébrica (\omterm{def:g10:algebra:equation}{equações}) e a funcional (\omterm{def:g10:reffunc:affine}{funções afins} do
\cref{ch:g10:reffunc}) --- e ensina as duas maneiras usuais de resolver um
\omterm{def:g10:lines:system}{sistema linear}.

\section{A equação de uma reta}

\begin{theorem}[Equação reduzida]\label{thm:g10:lines:reduced}
Em um \omterm{def:g10:coordgeom:system}{sistema de coordenadas}:
\begin{enumerate}
  \item toda reta não vertical tem uma \omterm{def:g10:algebra:equation}{equação} da forma
        \[
        y = mx + p ,
        \]
        com $m$ (o \emph{coeficiente angular}\index{coeficiente angular})
        e $p$ (o \emph{\omterm{def:g10:reffunc:affine}{coeficiente linear}}) unicamente determinados ---
        essa é a \emph{equação reduzida}\index{equação reduzida} da reta;
  \item toda reta vertical tem uma \omterm{def:g10:algebra:equation}{equação} $x = c$;
  \item um ponto está na reta exatamente quando suas \omterm{def:g10:coordgeom:system}{coordenadas}
        satisfazem a \omterm{def:g10:algebra:equation}{equação}.
\end{enumerate}
\end{theorem}
\admitted

\begin{proposition}[Coeficiente angular por dois pontos]\label{prop:g10:lines:slope}
O \omterm{def:g10:reffunc:affine}{coeficiente angular} da reta (não vertical) que passa por
$A(x_A, y_A)$ e $B(x_B, y_B)$, com $x_A \neq x_B$, é
\[
m = \frac{y_B - y_A}{x_B - x_A}
\qquad
\text{(``variação de $y$ sobre variação de $x$'').}
\]
\end{proposition}

\begin{proof}
Os dois pontos satisfazem $y = mx + p$, logo
$y_B - y_A = (mx_B + p) - (mx_A + p) = m(x_B - x_A)$; divida por
$x_B - x_A \neq 0$.
\end{proof}

\begin{omfigure}
\begin{tikzpicture}
\begin{axis}[omaxis,
  width=8.2cm, height=5.8cm,
  xmin=-1.4, xmax=4.4, ymin=-1.4, ymax=4.4,
  xtick={1,2,3}, ytick={1,3}]
  \addplot[omDef, thick, domain=-1.1:4.1] {x + 0.5};
  \addplot[omThm, thick, domain=-1.1:4.1] {3};
  \draw[omMeth, thick] (axis cs:2,-1.2) -- (axis cs:2,4.2);
  \node[omDef, font=\small, rotate=38] at (axis cs:2.7,3.6)
    {$y = x + 0.5$};
  \node[omThm, font=\small] at (axis cs:-0.55,3.3) {$y = 3$};
  \node[omMeth, font=\small, right] at (axis cs:2,-0.8) {$x = 2$};
\end{axis}
\end{tikzpicture}

{\small Três tipos de reta: inclinada ($m \neq 0$), horizontal ($m = 0$)
e vertical (sem \omterm{thm:g10:lines:reduced}{equação reduzida} --- sua \omterm{def:g10:algebra:equation}{equação} é $x = c$).}
\end{omfigure}

\begin{method}[Achar a equação da reta por dois pontos]\label{met:g10:lines:through}
Dados $A$ e $B$ com $x_A \neq x_B$:
\begin{enumerate}
  \item calcule o \omterm{def:g10:reffunc:affine}{coeficiente angular}
        $m = \dfrac{y_B - y_A}{x_B - x_A}$;
  \item escreva $y = mx + p$ e substitua as \omterm{def:g10:coordgeom:system}{coordenadas} de $A$ (ou de
        $B$) para achar $p$;
  \item confira a \omterm{def:g10:algebra:equation}{equação} final com o \emph{outro} ponto.
\end{enumerate}
Se $x_A = x_B$, a reta é vertical: sua \omterm{def:g10:algebra:equation}{equação} é $x = x_A$.
\end{method}

\begin{example}\label{ex:g10:lines:through}
Reta por $A(1, 5)$ e $B(3, 1)$. \omterm{def:g10:reffunc:affine}{Coeficiente angular}:
$m = \frac{1 - 5}{3 - 1} = \frac{-4}{2} = -2$. Então $y = -2x + p$; o
ponto $A$ dá $5 = -2 \times 1 + p$, logo $p = 7$. \omterm{def:g10:algebra:equation}{Equação}:
$y = -2x + 7$. Confira com $B$: $-2 \times 3 + 7 = 1$.
\end{example}

\begin{proposition}[Retas paralelas]\label{prop:g10:lines:parallel}
Duas retas não verticais são paralelas se, e somente se, têm o mesmo
\omterm{def:g10:reffunc:affine}{coeficiente angular}. Duas retas paralelas distintas não têm ponto comum;
duas retas de \omterm{def:g10:reffunc:affine}{coeficientes angulares} diferentes têm exatamente um.
\end{proposition}

\begin{proof}
Sejam $y = mx + p$ e $y = m'x + p'$ as duas retas. Um ponto comum
resolve $mx + p = m'x + p'$, isto é, $(m - m')x = p' - p$. Se
$m \neq m'$, isso tem exatamente uma solução
$x = \frac{p' - p}{m - m'}$: um ponto de \omterm{def:g10:numbers:interunion}{interseção}. Se $m = m'$, a
\omterm{def:g10:algebra:equation}{equação} fica $0 = p' - p$: nenhuma solução quando $p \neq p'$ (paralelas
distintas), e todo $x$ quando $p = p'$ (a mesma reta). O paralelismo (não
ter \omterm{def:g10:numbers:interunion}{interseção}, ou coincidir) é, portanto, equivalente a $m = m'$.
\end{proof}

\begin{remark}[Vetor diretor]
A reta $y = mx + p$ sobe $m$ unidades a cada unidade para a direita, de
modo que o \omterm{def:g10:vectors:vector}{vetor} $\vec u\,(1, m)$ é um \emph{vetor diretor}\index{vetor
diretor} da reta: a reta é paralela a $\vec u$. Duas retas são paralelas
exatamente quando seus \omterm{def:g10:vectors:vector}{vetores} diretores são \omterm{def:g10:vectors:collinear}{colineares}
(\cref{thm:g10:vectors:det}).
\end{remark}

\section{Sistemas lineares}

\begin{definition}[Sistema linear]\label{def:g10:lines:system}
Um \emph{sistema de duas equações lineares}\index{sistema linear} nas
incógnitas $x$ e $y$ é um par de \omterm{def:g10:algebra:equation}{equações}
\[
\begin{cases}
a x + b y = c \\
a' x + b' y = c'
\end{cases}
\]
a serem satisfeitas \emph{simultaneamente}. Uma \emph{solução} é um par
$(x, y)$ que satisfaz as duas. Cada \omterm{def:g10:algebra:equation}{equação} descreve uma reta, de modo
que resolver o sistema é encontrar a \omterm{def:g10:numbers:interunion}{interseção} de duas retas.
\end{definition}

\begin{method}[Substituição]\label{met:g10:lines:substitution}
\begin{enumerate}
  \item Isole uma incógnita em uma das \omterm{def:g10:algebra:equation}{equações} --- escolha a mais
        fácil, por exemplo exprima $y$ em \omterm{def:g10:functions:function}{função} de $x$;
  \item substitua essa expressão na \emph{outra} \omterm{def:g10:algebra:equation}{equação}, que passa a
        conter apenas $x$;
  \item resolva em $x$ e calcule em seguida $y$ pela etapa 1;
  \item confira o par nas duas \omterm{def:g10:algebra:equation}{equações} originais.
\end{enumerate}
\end{method}

\begin{example}\label{ex:g10:lines:substitution}
Resolva $\begin{cases} x + 2y = 8 \\ 3x - y = 3 \end{cases}$ por
substituição. A segunda \omterm{def:g10:algebra:equation}{equação} dá $y = 3x - 3$. Substituindo na
primeira:
\[
x + 2(3x - 3) = 8
\ \Longleftrightarrow\
x + 6x - 6 = 8
\ \Longleftrightarrow\
7x = 14
\ \Longleftrightarrow\
x = 2 ,
\]
e então $y = 3 \times 2 - 3 = 3$. A solução é o par $(2, 3)$.
Verificação: $2 + 2 \times 3 = 8$ e $3 \times 2 - 3 = 3$.
\end{example}

\begin{method}[Eliminação]\label{met:g10:lines:elimination}
\begin{enumerate}
  \item Multiplique cada \omterm{def:g10:algebra:equation}{equação} por um número bem escolhido, de modo que
        uma das incógnitas fique com coeficientes opostos nas duas
        \omterm{def:g10:algebra:equation}{equações};
  \item some as \omterm{def:g10:algebra:equation}{equações}: essa incógnita desaparece;
  \item resolva a \omterm{def:g10:algebra:equation}{equação} restante, de uma só incógnita, e substitua de
        volta para achar a outra;
  \item confira o par nas duas \omterm{def:g10:algebra:equation}{equações} originais.
\end{enumerate}
\end{method}

\begin{example}\label{ex:g10:lines:elimination}
Resolva $\begin{cases} 2x + 3y = 7 \\ 5x - 2y = 8 \end{cases}$ por
eliminação. Para eliminar $y$, multiplique a primeira \omterm{def:g10:algebra:equation}{equação} por $2$ e a
segunda por $3$:
\[
\begin{cases}
4x + 6y = 14 \\
15x - 6y = 24
\end{cases}
\]
Somando-as: $19x = 38$, logo $x = 2$. Substituindo em $2x + 3y = 7$:
$4 + 3y = 7$, logo $y = 1$. Solução: $(2, 1)$. Verificação na segunda
\omterm{def:g10:algebra:equation}{equação}: $5 \times 2 - 2 \times 1 = 8$.
\end{example}

\begin{omfigure}
\begin{tikzpicture}
\begin{axis}[omaxis,
  width=8.2cm, height=6.0cm,
  xmin=-0.9, xmax=4.9, ymin=-1.4, ymax=4.9,
  xtick={2}, ytick={1}]
  \addplot[omDef, thick, domain=-0.6:4.6] {(7 - 2*x)/3};
  \addplot[omThm, thick, domain=-0.5:2.9] {(5*x - 8)/2};
  \fill (axis cs:2,1) circle (2pt);
  \node[above right, font=\small] at (axis cs:2,1) {$(2, 1)$};
  \node[omDef, font=\small] at (axis cs:4.05,0.35) {$2x + 3y = 7$};
  \node[omThm, font=\small] at (axis cs:1.3,3.3) {$5x - 2y = 8$};
\end{axis}
\end{tikzpicture}

{\small Um \omterm{def:g10:lines:system}{sistema linear} visto geometricamente: cada \omterm{def:g10:algebra:equation}{equação} é uma reta,
e a solução $(2,1)$ do \cref{ex:g10:lines:elimination} é o seu ponto de
\omterm{def:g10:numbers:interunion}{interseção}.}
\end{omfigure}

\begin{remark}[Quantas soluções?]
Como duas retas, um \omterm{def:g10:lines:system}{sistema linear} tem em geral exatamente uma solução
(retas não paralelas) e, excepcionalmente, nenhuma (paralelas distintas)
ou uma infinidade (duas vezes a mesma reta). O critério do
\cref{thm:g10:vectors:det} decide: o sistema tem solução única
exatamente quando $ab' - a'b \neq 0$.
\end{remark}

\begin{example}[Modelar com um sistema]\label{ex:g10:lines:modeling}
Três cadernos e duas canetas custam $8$; um caderno e quatro canetas
custam $6$. Seja $x$ o preço de um caderno e $y$ o de uma caneta:
\[
\begin{cases}
3x + 2y = 8 \\
x + 4y = 6 .
\end{cases}
\]
Da segunda \omterm{def:g10:algebra:equation}{equação}, $x = 6 - 4y$; substituindo, $3(6 - 4y) + 2y = 8$,
logo $18 - 10y = 8$, o que dá $y = 1$ e depois $x = 2$. Um caderno custa
$2$ e uma caneta custa $1$.
\end{example}

\section{Exercícios}

\begin{exercise}[$\star$]\label{exo:g10:lines:1}
Para cada reta, leia o \omterm{def:g10:reffunc:affine}{coeficiente angular} e o \omterm{def:g10:reffunc:affine}{coeficiente linear} e
esboce-a:
\[
y = 2x - 3, \qquad
y = -\tfrac13 x + 1, \qquad
y = 4, \qquad
x = -2 .
\]
\end{exercise}

\begin{exercise}[$\star$]\label{exo:g10:lines:2}
O ponto $A(2, 7)$ pertence à reta $y = 3x + 1$? E $B(-1, -1)$? E o ponto
$C(4, 13)$?
\end{exercise}

\begin{exercise}[$\star$]\label{exo:g10:lines:3}
Encontre a \omterm{thm:g10:lines:reduced}{equação reduzida} da reta que passa por:
\[
\text{(a) } A(0, 2) \text{ e } B(3, 8);
\qquad
\text{(b) } A(-1, 5) \text{ e } B(2, -4);
\qquad
\text{(c) } A(4, 1) \text{ e } B(4, 6).
\]
\end{exercise}

\begin{exercise}[$\star$]\label{exo:g10:lines:4}
Entre as retas $y = 3x - 1$, $y = -3x + 2$, $y = 3x + 5$ e
$y = \frac{6x + 2}{2}$, quais são paralelas entre si? Quais são, na
verdade, iguais?
\end{exercise}

\begin{exercise}[$\star$]\label{exo:g10:lines:5}
Resolva por substituição:
\[
\begin{cases}
y = 2x - 1 \\
3x + y = 9
\end{cases}
\qquad\text{e depois}\qquad
\begin{cases}
x - y = 4 \\
2x + 3y = 3 .
\end{cases}
\]
\end{exercise}

\begin{exercise}[$\star$]\label{exo:g10:lines:6}
Resolva por eliminação:
\[
\begin{cases}
x + y = 10 \\
x - y = 4
\end{cases}
\qquad\text{e depois}\qquad
\begin{cases}
3x + 2y = 12 \\
2x + 5y = 8 .
\end{cases}
\]
\end{exercise}

\begin{exercise}[$\star\star$]\label{exo:g10:lines:7}
Determine, sem resolvê-los, quantas soluções cada sistema tem:
\[
\begin{cases}
2x - y = 3 \\
-4x + 2y = -6
\end{cases}
\qquad
\begin{cases}
2x - y = 3 \\
-4x + 2y = 1
\end{cases}
\qquad
\begin{cases}
2x - y = 3 \\
x + y = 0 .
\end{cases}
\]
\end{exercise}

\begin{exercise}[$\star\star$]\label{exo:g10:lines:8}
Foram vendidos duzentos ingressos para uma peça escolar, uns a $5$
(crianças) e outros a $9$ (adultos), totalizando $1352$. Quantos
ingressos de cada tipo foram vendidos?
\end{exercise}

\begin{exercise}[$\star\star$]\label{exo:g10:lines:9}
Seja $d$ a reta $y = 2x - 3$ e seja $A(1, 4)$.
\begin{enumerate}
  \item Dê a \omterm{def:g10:algebra:equation}{equação} da reta $d'$ paralela a $d$ que passa por $A$.
  \item Calcule o ponto de \omterm{def:g10:numbers:interunion}{interseção} de $d'$ com o eixo $x$.
\end{enumerate}
\end{exercise}

\begin{exercise}[$\star\star\star$]\label{exo:g10:lines:10}
Sejam $A(0, 0)$, $B(6, 2)$ e $C(2, 4)$.
\begin{enumerate}
  \item Encontre as \omterm{def:g10:algebra:equation}{equações} das medianas do triângulo $ABC$ que partem
        de $A$ e de $B$ (uma mediana liga um vértice ao \omterm{prop:g10:coordgeom:midpoint}{ponto médio} do
        lado oposto).
  \item Calcule seu ponto de \omterm{def:g10:numbers:interunion}{interseção} $G$ e verifique que $G$ também
        está na terceira mediana. (Você deve encontrar que as \omterm{def:g10:coordgeom:system}{coordenadas}
        de $G$ são as médias das de $A$, $B$, $C$.)
\end{enumerate}
\end{exercise}

\section{Problema: Uma solução, nenhuma ou uma infinidade}

\begin{problem}[{Problema de fim de semana --- duas retas têm três
destinos possíveis, um determinante decide entre eles, e o mais
antigo livro-texto do mundo já sabia disso}]
\label{pb:g10:lines:1}
Duas \omterm{def:g10:algebra:equation}{equações} lineares, duas incógnitas: o par de retas que elas
desenham pode se cruzar uma vez, nunca se cruzar, ou ser uma só e
mesma reta --- e \emph{todo} \omterm{def:g10:lines:system}{sistema linear} herda um desses três
destinos. Este problema os classifica, apresenta o número que
decide (um velho conhecido do capítulo dos \omterm{def:g10:vectors:vector}{vetores}), resolve os
faisões e coelhos, com dois mil anos, dos \emph{Nove Capítulos}
chineses, e termina com os vértices de regiões --- o primeiro
passo da otimização usada hoje por toda companhia aérea e toda
fábrica.

\medskip
\noindent\textbf{Parte I --- A reta, com fluência.}
\begin{enumerate}
  \item Encontre a \omterm{thm:g10:lines:reduced}{equação reduzida} da reta que passa por
        $(1, 2)$ e $(3, 8)$
        (\cref{met:g10:lines:through}).
  \item Entre $y = 3x - 1$, $y = 3x + 4$ e $y = -x + 7$: quais
        duas retas são paralelas
        (\cref{prop:g10:lines:parallel}) e onde se encontram os
        pares não paralelos? (Calcule uma \omterm{def:g10:numbers:interunion}{interseção}.)
  \item A reta que passa por $(2, 1)$ e $(2, 5)$: por que ela não
        tem \omterm{thm:g10:lines:reduced}{equação reduzida} $y = mx + p$ e qual é a sua \omterm{def:g10:algebra:equation}{equação}?
  \item O ponto $(4, 11)$ está na reta $y = 3x - 1$? E
        $(-1, -5)$?
  \item A reta que passa por $(0, 3)$ com \omterm{def:g10:reffunc:affine}{coeficiente angular}
        $-2$: dê sua \omterm{def:g10:algebra:equation}{equação}, seus dois pontos de corte com os
        eixos e a área do triângulo que ela recorta no primeiro
        quadrante.
\end{enumerate}

\medskip
\noindent\textbf{Parte II --- Três destinos.}
\begin{enumerate}[resume]
  \item Resolva por substituição
        (\cref{met:g10:lines:substitution}):
        $\begin{cases} y = 2x - 3 \\ 3x + y = 7 \end{cases}$
  \item Resolva por eliminação
        (\cref{met:g10:lines:elimination}):
        $\begin{cases} 2x + 3y = 8 \\ 5x - 3y = -1 \end{cases}$
  \item Classifique --- e interprete com retas --- os dois
        sistemas
        \[
        \begin{cases} 2x - y = 3 \\ 4x - 2y = 6 \end{cases}
        \qquad\text{e}\qquad
        \begin{cases} 2x - y = 3 \\ 4x - 2y = 1 . \end{cases}
        \]
        Enuncie a tricotomia completa: concorrentes, paralelas,
        idênticas --- uma solução, nenhuma, uma infinidade.
  \item Para o sistema geral $ax + by = e$, $cx + dy = f$, o
        número que decide é $ad - bc$ --- o mesmo determinante do
        teste de colinearidade do \cref{pb:g10:vectors:1}.
        Explique a coincidência (quais dois \omterm{def:g10:vectors:vector}{vetores} são
        \omterm{def:g10:vectors:collinear}{colineares} exatamente quando as retas são paralelas ou
        idênticas?) e calcule $ad - bc$ para os três sistemas das
        questões 7 e 8.
  \item Para que valor de $m$ o sistema
        $\begin{cases} x + my = 3 \\ 2x + 6y = 7 \end{cases}$
        não tem solução? Resolva o sistema para todos os demais
        valores de $m$\,\dots\ ou, ao menos, explique como você
        faria.
\end{enumerate}

\medskip
\noindent\textbf{Parte III --- Os Nove Capítulos.}
\begin{enumerate}[resume]
  \item Dos \emph{Nove Capítulos sobre a Arte Matemática}
        (China, há cerca de 2000 anos): uma gaiola guarda faisões
        e coelhos --- $35$ cabeças e $94$ patas ao todo. Quantos
        de cada?
  \item Uma barraca de sucos mistura um xarope com $60\,\%$ de
        fruta e uma bebida com $20\,\%$ de fruta para obter
        $30$ L a $35\,\%$ de fruta. Quantos litros de cada?
  \item Um número de dois algarismos tem soma de algarismos $12$;
        trocar seus algarismos de posição o diminui em $18$.
        Encontre-o, usando a álgebra dos algarismos do problema
        de fim de semana correspondente, no volume do ensino
        fundamental, para montar o sistema.
  \item Na padaria, $3$ cafés e $2$ croissants custam $9.10$
        euros; $2$ cafés e $3$ croissants custam $8.90$ euros.
        Resolva com elegância: o que \emph{somar} as duas
        \omterm{def:g10:algebra:equation}{equações} e o que \emph{subtraí-las} dizem de imediato?
  \item Uma banca de feira anuncia: $2$ maçãs $+ 1$ pera por
        $5$ euros e $4$ maçãs $+ 2$ peras por $9$ euros. Resolva
        --- ou melhor, explique o que o destino do sistema revela
        sobre a aritmética da banca.
\end{enumerate}

\medskip
\noindent\textbf{Parte IV --- Os vértices mandam.}
\begin{enumerate}[resume]
  \item A reta $y = 3x - 1$ parte o plano em dois. De que lado
        está a origem? Descreva o conjunto $y > 3x - 1$ e como se
        testa um ponto contra ele.
  \item Desenhe a região definida pelas quatro restrições
        \[
        x \geq 0, \qquad y \geq 0, \qquad x + y \leq 4,
        \qquad y \leq 2x + 1,
        \]
        e calcule seus quatro vértices.
  \item O lucro de uma oficina é $P = 3x + 2y$ na região da
        questão 17. Calcule $P$ nos quatro vértices. Admitindo
        que o \omterm{def:g10:functions:extrema}{máximo} de uma grandeza linear como essa sobre uma
        região desse tipo é sempre atingido em um vértice
        (imagine as retas $3x + 2y = $ constante varrendo o
        plano), dê o plano de produção ótimo.
  \item Os dois táxis do problema de fim de semana sobre os dois
        termômetros, no volume do ensino fundamental, como
        sistema: escreva e resolva
        $\begin{cases} y = 1.5x + 3 \\ y = 2x \end{cases}$ e
        interprete a solução.
  \item Final, o dicionário: uma \omterm{def:g10:algebra:equation}{equação} $\leftrightarrow$ uma
        reta; um sistema $\leftrightarrow$ duas retas com três
        destinos, decididos por $ad - bc$; muitas inequações
        $\leftrightarrow$ uma região cujos vértices carregam o
        ótimo. Enuncie isso em três frases --- você acaba de
        percorrer o saguão de entrada da \emph{álgebra linear} e
        da \emph{programação linear}, ambas construídas por
        inteiro nos volumes de graduação.
\end{enumerate}
\end{problem}
